\section{Planejamento da avaliação}

\subsection{Seleção e priorização}

A avaliação adota o modelo de qualidade de produto da ISO/IEC
25010:2011 e utiliza a ISO/IEC 25023:2016 como referência para a
definição posterior das medidas \cite{iso25010_2011,iso25023_2016}.
Foram selecionadas Eficiência de desempenho, Portabilidade e
Usabilidade, desdobradas em comportamento temporal, utilização de
recursos, adaptabilidade, acessibilidade e apreensibilidade.
Como hipótese de trabalho, considera-se o acesso por estudantes em
computadores de uma escola pública, sujeitos a possíveis limitações de
processamento, memória e conectividade. O cenário orienta a seleção e a
profundidade dos ensaios, mas não é apresentado como resultado de uma
observação de campo nem como representação de todas as escolas.

A priorização utiliza uma matriz com notas de 1 a 3 e três critérios:
impacto no estudante, com peso 3; risco no contexto, com peso 2; e
viabilidade de evidência, com peso 1. A pontuação
$P=3I+2R+V$ classifica as subcaracterísticas em prioridade alta
($P\geq15$), média ($12\leq P\leq14$) ou baixa ($P\leq11$).
Comportamento temporal obteve 18 pontos, utilização de recursos 17,
adaptabilidade e acessibilidade 15, e apreensibilidade 12. A
ponderação privilegia impacto e risco sem permitir que a facilidade
de medição determine isoladamente o escopo.

O principal trade-off é entre cobertura e profundidade. Em vez de
executar todas as técnicas em todas as combinações de sistema e
navegador, cada subcaracterística recebe ambientes e procedimentos
coerentes com a evidência necessária.

\subsection{Escopo e profundidade}

O objeto é a plataforma web da Khan Academy acessada por navegadores
em computadores. Aplicativos nativos, navegadores móveis, painéis de
pais, professores e administradores, infraestrutura interna, testes
ativos de segurança e avaliação pedagógica do conteúdo permanecem
fora do escopo. O recorte concentra-se em localizar conteúdo,
reproduzir uma videoaula, responder a um exercício e consultar o
progresso da própria conta.

A matriz de ambientes contempla Chrome no Windows, Firefox no
Windows, Safari no macOS e Firefox no Ubuntu. Essa seleção cobre os
motores Blink, Gecko e WebKit. Os ensaios de acessibilidade utilizam
NVDA, VoiceOver e Orca nas combinações compatíveis, com navegação por
teclado e inspeção automatizada apenas como apoio.

Comportamento temporal e utilização de recursos serão medidos em um
ambiente de referência, com uma execução de aquecimento e cinco
repetições válidas por tarefa e condição de rede. A comparação entre a
conexão sem limitação artificial e um perfil controlado de conexão
limitada operacionaliza a preocupação do cenário escolar com barreiras
de conectividade, sem extrapolar os resultados para condições não
avaliadas. Adaptabilidade será verificada nos quatro ambientes e em duas
resoluções de computador.
A apreensibilidade terá caráter exploratório, com três a cinco
participantes adultos sem experiência anterior declarada com a
plataforma. Configurações, versões, dados brutos e ocorrências serão
registrados para permitir auditoria.

Os resultados serão válidos apenas para os fluxos, versões,
equipamentos, condições de rede e tecnologias assistivas
efetivamente utilizados. A amostra de apreensibilidade não representa
estatisticamente estudantes de escolas públicas, e os ensaios com
leitores de tela não comprovam conformidade integral com as WCAG.

\subsection{Relação com os Objetivos de Desenvolvimento Sustentável}

O estudo relaciona-se de forma contextual com o ODS 4, Meta 4.5, e
com o ODS 10, Meta 10.2 \cite{un_sdg4_5,un_sdg10_2}. Ensaios de
acessibilidade e apreensibilidade podem identificar barreiras
técnicas pertinentes ao acesso educacional, enquanto eficiência e
adaptabilidade podem revelar limitações em computadores e conexões
restritos.

Essa relação não significa que a avaliação mede os indicadores
oficiais. O indicador 4.5.1 trata de índices de paridade aplicados a
indicadores educacionais desagregáveis; o indicador 10.2.1 mede a
proporção de pessoas abaixo de 50\% da renda mediana, por sexo, idade
e deficiência. Portanto, os resultados caracterizam a qualidade
observada do software, sem demonstrar efeitos sobre participação
educacional, renda ou inclusão social.
